\clearpage\SourceLeaf{85}%
\pdfbookmark[1]{Calculation of Passover}{osterfest1802}
\begin{center}
{\large CALCULATION OF PASSOVER.\par}
\medskip\hrule\medskip
{\small Monatliche Correspondenz zur Beförderung der Erd- und Himmels-Kunde,\par
edited by Baron \textsc{von Zach}. May 1802.\par}
\medskip\hrule
\end{center}\medskip
The 15th of Nisan in the Jewish year $A$, on which Jews celebrate Passover, falls in the year $A-3760=B$ of the Christian era. The following purely arithmetical rule determines the corresponding day of the month.

Divide $12A+17$, or equivalently here $12B+12$, by 19, and call the remainder $a$. Then divide $A$ or $B$ by 4 and set the remainder $=b$. Calculate the value
\begin{center}\small
\begin{tabular}{lll}
\multicolumn{1}{c}{of} & \multicolumn{1}{c}{or of} & \shortstack{Decimal values\\as ordinary fractions}\\[5pt]
$32.0440932$ & $20.0955877$ & $32\frac{4343}{98496}\quad20\frac{9415}{98496}$\\[5pt]
$+1.5542418a$ & $+1.5542418a$ & $1\frac{272953}{492480}$\\[5pt]
$+0.25\quad b$ & $+0.25\quad b$ & $\frac14$\\[5pt]
$-0.003177794A$ & $-0.003177794B$ & $\frac{313}{98496}$
\end{tabular}
\end{center}
and set it $=M+m$, so that $M$ denotes the whole-number part and $m$ the (decimal) fractional part. Finally, divide $M+3A+5b+5$ or $M+3B+5b+1$ by 7 and set the remainder $=c$. The following four cases must now be distinguished.

I. If $c=2$, or if it is 4 or 6, Passover falls on the $M+1^{\mathrm{th}}$ of March \textit{Old Style}. Write this as the $M-30^{\mathrm{th}}$ of April when $M>30$ because of \textit{Adu}.
\SourceLeaf{86}%

II. If $c=1$, and if $a>6$ and $m\geqq0.63287037$
$(\frac{311676}{492480}=\frac{1367}{2160})$ both hold, Passover falls on the $M+2^{\mathrm{th}}$ of March \textit{O.S.} because of \textit{Getred}.

III. If $c=0$, and if $a>11$ and $m\geqq0.89772376$
$(\frac{442111}{492480}=\frac{23269}{25920})$ both hold, Passover falls on the $M+1^{\mathrm{th}}$ of March \textit{O.S.} because of \textit{Batu Thakpad}.

IV. In all other cases Passover falls on the $M^{\mathrm{th}}$ of March Old Style.

\begin{center}\rule{.25\linewidth}{.3pt}\end{center}

\textit{First remark.} These rules also determine the date of 1 \textit{Tisri}, or New Year’s Day, which always falls 163 days after Passover of the preceding year.

\textit{Second remark.} The year $A$ is a common year (of 12 months) if $a<12$, and a leap year (of 13 months) if $a>11$.

\Needspace{13\baselineskip}
\begin{center}\small
\setlength{\tabcolsep}{4pt}
\begin{tabular}{@{}l@{\hspace{8pt}}l@{}p{14pt}@{}l@{}}
\textit{Example of applying these rules:} & $A=5562$ & & $B=1802$\\
& $12A+17=66761$ & & $12B+12=21636$\\
\multicolumn{4}{c}{Dividing by 19 gives $a=14$}\\
\multicolumn{4}{c}{Dividing 5562 or 1802 by 4 gives $b=2$.}\\
Hence the values from the formula above & $32.0440932$ & & $20.0955877$\\
& $+21.7593852$ & & $21.7593852$\\
& $+\phantom{0}0.5$ & & $0.5$\\
& $-17.6748903$ & & $-\phantom{0}5.7263848$\\\cline{2-2}\cline{4-4}
& $36.6285881$ & & $36.6285881$
\end{tabular}
\end{center}
so $M=36,\quad m=0.6285881$,
\[
\begin{aligned}
M+3A+5b+5&=36+16686+10+5=16737\\
M+3B+5b+1&=36+\phantom{0}5406+10+1=\phantom{0}5453
\end{aligned}
\]
and division by 7 gives $c=0$.

Since $m$ is less than $0.89772376$ here, rule III cannot apply. By rule IV, then, Passover falls on the 36th of March Old Style or the 48th of March New Style, that is, the 17th of April.

In \textit{most} cases it suffices to calculate the formula above to only about 2 decimal places.

\begin{center}\rule{.25\linewidth}{.3pt}\end{center}
